\section{Causal structure and the birth-front normal form}
\label{sec:causal}

\subsection{Shortest traces have connected dependency graphs}

\begin{theorem}[Causal connectedness]
\label{thm:causal-connected}
Assume strong disjoint commutation.  If $\tau$ is a shortest first-unlocking
trace and $r$ is a terminal reduction exposed at its end, then
$H(\tau,r)$ is connected.
\end{theorem}

\begin{proof}
Suppose $H(\tau,r)$ is disconnected.  Let $C$ be a connected component not
containing the terminal vertex $\star$.  Every action in $C$ has footprint
disjoint from every action outside $C$ and from the terminal reduction.

Use the commutation diamond repeatedly to move every action of $C$ to the right,
preserving the relative order inside $C$ and the relative order of all actions
outside $C$.  Each adjacent swap is legal, preserves the reached state, and
preserves the relevant footprints.  We obtain a trace-equivalent word
\[
   \tau'\,\sigma_C,
\]
where $\sigma_C$ consists exactly of the actions in $C$.

The terminal reduction is legal after $\tau'\sigma_C$.  Since its footprint is
disjoint from every footprint in $\sigma_C$, terminal locality moves it left
across the entire block.  Hence it is already legal after $\tau'$.  Deleting
$\sigma_C$ gives a strictly shorter trace that exposes a terminal reduction;
if one was exposed at an earlier prefix of $\tau'$, truncate at its first such
prefix.  In either case we obtain a strictly shorter first-unlocking trace from
$x_0$, contradicting the minimality of $\tau$.
\end{proof}

\begin{corollary}
Every birth in a shortest trace has a path, through overlapping later action
footprints, to the terminal reduction.
\end{corollary}

This corollary is qualitative.  It rules out irrelevant disconnected work, but
it does not bound the number of births.  A connected dependency tree can have
many leaves.

\subsection{Why connected prefixes are false}

The following proposition records the proof-audit failure that led to the birth
parameter.

\begin{proposition}[Branching obstruction]
\label{prop:branching-obstruction}
There exists a strongly local reversible rewrite system with a shortest
first-unlocking trace of length three such that every first-unlocking trace of
length three has birth number two.  In particular, no commuting-equivalent
first-unlocking trace has connected chronological prefixes.
\end{proposition}

\begin{proof}
Use three Boolean sites $0,1,2$, initially all zero.  Action $A$ toggles site
$0$ and has footprint $\{0\}$.  Action $B$ toggles site $2$ and has footprint
$\{2\}$.  Action $C$ is legal exactly when sites $0$ and $2$ are one; it toggles
site $1$ and has footprint $\{0,1,2\}$.  The terminal reduction is legal exactly
when site $1$ is one.

The moves are involutions, and disjoint supports commute.  Neither $A$ nor $B$
alone enables $C$, so at least two moves are required before $C$.  The words
$ABC$ and $BAC$ unlock the reduction in three moves.  Their first two footprints
are disjoint, so each has two births.  No trace with one birth can set both sites
$0$ and $2$ before applying $C$.
\end{proof}

The fixture is implemented as \code{branching\_example()} in the prototype.  It
is not asserted to be a knot diagram; its purpose is to refute an invalid
general lemma before that lemma is imported into knot topology.

\subsection{Births commute to the front}

\begin{lemma}[Births are pairwise disjoint]
\label{lem:birth-disjoint}
The footprints of distinct births in one trace are disjoint.
\end{lemma}

\begin{proof}
A later birth is disjoint from the union of every earlier footprint, including
the footprint of any earlier birth.
\end{proof}

\begin{lemma}[Left commutation of a birth]
\label{lem:birth-left}
A birth action can be commuted left across every action preceding it.
\end{lemma}

\begin{proof}
By definition, its footprint is disjoint from each preceding footprint.  Apply
the backward-and-forward commutation diamond one adjacent predecessor at a
time.  Footprints remain unchanged under each swap, so the next swap remains
valid.
\end{proof}

\begin{theorem}[Birth-front normal form]
\label{thm:birth-front}
Let $\tau$ be a first-unlocking trace of length $k$ and birth number $b$ in a
strongly local reversible system.  There is a first-unlocking trace
$\widehat\tau$ of length at most $k$ and birth number at most $b$ such that:
\begin{enumerate}
\item all births of $\widehat\tau$ occur before every nonbirth action;
\item after the birth block, every action footprint intersects the union of all
      earlier footprints.
\end{enumerate}
\end{theorem}

\begin{proof}
Identify the births in the original trace.  By Lemma~\ref{lem:birth-left}, move
each birth left across all predecessors.  By
Lemma~\ref{lem:birth-disjoint}, births commute with one another, so choose any
fixed canonical order for them.  Preserve the relative order of all nonbirth
actions.

Consider an original nonbirth action $a_i$.  It overlapped at least one earlier
action.  If that earlier action was a birth, it has moved farther left.  If it
was a nonbirth, their relative order was preserved.  Thus $a_i$ still meets the
accumulated earlier footprint in the reordered word.  Strong commutation shows
that the reordered word is legal and reaches the same final state.

It is possible that the terminal reduction becomes legal at an earlier prefix
of the reordered word.  In that case truncate at the first such prefix.  The
length and number of births can only decrease.  The truncated word is
first-unlocking and has the two stated properties.
\end{proof}

\begin{remark}[Normal form versus a chronological invariant]
The birth number is defined on a trace, not just its final state.  The theorem
shows that it has algorithmic meaning: it is an upper bound on the number of
global choices needed before all remaining choices are local.  One may minimize
$b$ over shortest traces to obtain an input parameter, but the search does not
need to know that minimum in advance; it can deepen in $b$.
\end{remark}

\subsection{Accumulated support strictly improves last-touch search}

The current ProveIt recursion follows the neighborhood of the immediately
previous RIII move.  The next proposition shows that even the one-birth normal
form is strictly more general in local rewrite systems.

\begin{proposition}[Strict separation]
\label{prop:last-vs-accumulated}
There is a strongly local reversible system with a one-birth first-unlocking
trace of length three, but no first-unlocking trace of length at most three in
which consecutive action footprints always intersect.
\end{proposition}

\begin{proof}
Use sites $0,1,2,3$, initially zero.  Action $A$ toggles sites $0$ and $1$ and
has footprint $\{0,1\}$.  Action $B$ requires site $0=1$, toggles site $2$, and
has footprint $\{0,2\}$.  Action $C$ requires site $1=1$, toggles site $3$, and
has footprint $\{1,3\}$.  The terminal reduction requires sites $2=3=1$.

After $A$, both $B$ and $C$ are legal.  They commute because their footprints
are disjoint.  Either $ABC$ or $ACB$ unlocks the terminal reduction.  There is
one birth, namely $A$, and each leaf meets the accumulated footprint of $A$.
However, $B$ and $C$ do not meet one another, so neither length-three ordering
has pairwise-intersecting consecutive footprints.  No shorter word can set both
sites $2$ and $3$.
\end{proof}

The fixture is \code{accumulated\_strict\_example()} in the prototype.  For knot
diagrams, this proposition motivates the accumulated-footprint search; it does
not by itself exhibit a concrete RIII diagram with the same separation.
